% Fig. 3 划分方差 vs 种子方差（a）；单一划分上的“LLM 领先”在重复划分上消失（b）
% 数据：data/fig3a_variance.dat、data/fig3b_cic_lgbm.dat（由 prep_figs.py 生成）；编译：pdflatex fig3_split_variance.tex
\documentclass[10pt]{standalone}
% 全部论文图共用的样式（颜色固定跟随来源；字号按 elsarticle 正文 10pt 设定）
% 颜色来自经过色觉校验的分类调色板，顺序固定：
%   类别权重 = 1 蓝，ROS = 2 橙，SMOTE = 3 青，Qwen = 4 黄，Llama = 5 品红；不增强 = 中性灰
\usepackage{pgfplots}
\usepgfplotslibrary{groupplots}
\pgfplotsset{compat=1.18}
\usetikzlibrary{plotmarks,calc,patterns,arrows.meta,positioning,shapes.geometric,fit,backgrounds}
\definecolor{cCW}{HTML}{2A78D6}
\definecolor{cROS}{HTML}{EB6834}
\definecolor{cSMOTE}{HTML}{1BAF7A}
\definecolor{cQWEN}{HTML}{EDA100}
\definecolor{cLLAMA}{HTML}{E87BA4}
\definecolor{cNONE}{HTML}{52514E}
\definecolor{cGRID}{HTML}{E4E3DF}
\definecolor{cAXIS}{HTML}{8A8984}
\definecolor{cTEXT}{HTML}{2B2A28}
\pgfplotsset{
  paper/.style={
    font=\footnotesize, color=cTEXT,
    axis line style={cAXIS, line width=0.4pt},
    tick style={cAXIS, line width=0.4pt},
    major tick length=2pt,
    grid style={cGRID, line width=0.4pt},
    ymajorgrids,
    tick label style={font=\scriptsize},
    label style={font=\footnotesize},
    title style={font=\footnotesize\bfseries, yshift=-2pt},
    legend style={font=\scriptsize, draw=none, fill=none, /tikz/every even column/.append style={column sep=8pt}},
    legend cell align=left,
  },
  errbar/.style={error bars/.cd, y dir=both, y explicit,
    error bar style={line width=0.5pt, cTEXT!75}, error mark options={rotate=90, mark size=1.4pt, line width=0.5pt, cTEXT!75}},
}

\begin{document}
\begin{tikzpicture}
\begin{groupplot}[
  paper,
  group style={group size=2 by 1, horizontal sep=1.6cm},
  height=6.0cm,
]
% ---------------- (a) 种子间 sd 与划分间 sd（哑铃图）
\nextgroupplot[
  width=6.4cm,
  title={(a) Seed vs.\ split variability},
  xmin=0, xmax=0.05, xtick={0,0.01,0.02,0.03,0.04,0.05}, scaled x ticks=false,
  xticklabel style={/pgf/number format/fixed, /pgf/number format/fixed zerofill, /pgf/number format/precision=2},
  xlabel={SD of rare-class macro-F1},
  ymajorgrids=false, xmajorgrids,
  ymin=-0.7, ymax=11.7, ytick={0,...,11},
  yticklabels={CIC · XGB, CIC · LGBM, CIC · RF-grid, CIC · SAORF,
               NF · XGB, NF · LGBM, NF · RF-grid, NF · SAORF,
               CSE · XGB, CSE · LGBM, CSE · RF-grid, CSE · SAORF},
  y tick label style={font=\scriptsize},
  legend style={at={(0.5,-0.2)}, anchor=north},
  legend columns=1,
]
  % 每行一条连接线（种子间 sd → 划分间 sd）
  \pgfplotstableread{data/fig3a_variance.dat}\vartab
  \pgfplotsforeachungrouped \i in {0,...,11} {
    \pgfplotstablegetelem{\i}{y}\of\vartab \let\yy\pgfplotsretval
    \pgfplotstablegetelem{\i}{seed_sd}\of\vartab \let\xa\pgfplotsretval
    \pgfplotstablegetelem{\i}{split_sd}\of\vartab \let\xb\pgfplotsretval
    \edef\tmp{\noexpand\draw[cAXIS, line width=0.9pt] (axis cs:\xa,\yy) -- (axis cs:\xb,\yy);}\tmp
  }
  \addplot[only marks, mark=o, mark size=2.3pt, mark options={cNONE, line width=0.7pt, fill=white}] table[x=seed_sd, y=y]{data/fig3a_variance.dat};
  \addplot[only marks, mark=*, mark size=2.3pt, mark options={cCW, draw=white, line width=0.4pt}] table[x=split_sd, y=y]{data/fig3a_variance.dat};
  \legend{Across 10 seeds (one split), Across 5 splits}
  \draw[cGRID, line width=0.5pt] (axis cs:0,3.5) -- (axis cs:0.05,3.5);
  \draw[cGRID, line width=0.5pt] (axis cs:0,7.5) -- (axis cs:0.05,7.5);
% ---------------- (b) CIC-IDS2017 · LightGBM：原划分 vs 重复划分
\nextgroupplot[
  width=7.6cm,
  title={(b) CIC-IDS2017 $\cdot$ LightGBM},
  xmin=-0.5, xmax=3.5, xtick={0,1,2,3},
  xticklabels={Original split, Split 1, Split 2, Split 3},
  ymin=0.565, ymax=0.725, ytick={0.58,0.60,...,0.72},
  yticklabel style={/pgf/number format/fixed, /pgf/number format/fixed zerofill, /pgf/number format/precision=2},
  ylabel={Rare-class macro-F1}, ylabel style={yshift=-3pt},
  legend style={at={(0.5,-0.2)}, anchor=north, /tikz/every even column/.append style={column sep=6pt}},
  legend columns=3,
  error bars/error bar style={line width=0.5pt},
]
  \fill[cGRID!55] (axis cs:-0.5,0.565) rectangle (axis cs:0.5,0.725);
  \addplot[only marks, mark=square*, mark size=1.9pt, mark options={cNONE, draw=white, line width=0.3pt},
           x filter/.expression={x-0.28}, errbar] table[x=x, y=none, y error=none_sd]{data/fig3b_cic_lgbm.dat};
  \addplot[only marks, mark=*, mark size=2.2pt, mark options={cCW, draw=white, line width=0.3pt},
           x filter/.expression={x-0.14}, errbar] table[x=x, y=cw, y error=cw_sd]{data/fig3b_cic_lgbm.dat};
  \addplot[only marks, mark=diamond*, mark size=2.7pt, mark options={cROS, draw=white, line width=0.3pt},
           x filter/.expression={x}, errbar] table[x=x, y=ros, y error=ros_sd]{data/fig3b_cic_lgbm.dat};
  \addplot[only marks, mark=triangle*, mark size=2.6pt, mark options={cQWEN, draw=cTEXT!60, line width=0.3pt},
           x filter/.expression={x+0.14}, errbar] table[x=x, y=qwen, y error=qwen_sd]{data/fig3b_cic_lgbm.dat};
  \addplot[only marks, mark=pentagon*, mark size=2.6pt, mark options={cLLAMA, draw=cTEXT!60, line width=0.3pt},
           x filter/.expression={x+0.28}, errbar] table[x=x, y=llama, y error=llama_sd]{data/fig3b_cic_lgbm.dat};
  \legend{No augmentation, Class weight, ROS, LLM (Qwen2.5-7B), LLM (Llama-3.1-8B)}
\end{groupplot}
\end{tikzpicture}
\end{document}
